\begin{textsample}{NEG Dim 4 – pb – Score -3.00 – pb/pb\_ef\_7.txt}  \label{ex:f4_neg_123}
5.1.5 Avaliação A avaliação deve ser pensada como uma oportunidade de auxílio para a promoção do desenvolvimento, mas também como um recurso de acompanhamento progressivo desse desenvolvimento na formação de um sujeito em sua totalidade, enquanto ser humano, enquanto ser atuante e participativo em espaços diversos de interação.
Para tanto, é importante que o professor perceba a avaliação como um ato dialógico e processual, indicador de caminhos norteadores de possibilidades que irão contribuir para o crescimento do aprendente.
Deve estar, portanto, associada, diretamente, ao planejamento, numa lógica que segue a sequência de observação, diagnóstico e, por fim, tomada de decisões para atuação, buscando o desenvolvimento de ações estratégicas e/ou reorientações da prática pedagógica que promovam progressos a partir do suprimento das necessidades observadas.
Em Língua Portuguesa, essa perspectiva progressiva e formativa da avaliação deve permear os eixos da leitura, oralidade, escrita e análise linguística/semiótica.
A avaliação no eixo da leitura deve ser pensada no crescimento \textbf{gradual} da análise da compreensão leitora voltada à construção dos sentidos dos textos.
Isso precisa acontecer no dia a dia da sala de aula, partindo da oferta de gêneros textuais diversificados e de atividades que explorem, por exemplo, aspectos inferenciais, semânticos e coesivos nos textos.
É indispensável que tais atividades sejam pensadas a partir de instruções que explorem relevância, clareza, argumentação e síntese em níveis diferentes, conduzindo o aluno a atender às solicitações de forma exitosa e dando condições, ao professor, de \textbf{mensurar} a capacidade leitora desses educandos.
Vale ressaltar que não se trata de propor somente atividades puramente didáticas, mas também de promover momentos de partilha de leituras para a fruição, o despertar para o gosto pela leitura por prazer, pelo riso despretensioso, pela apreciação, para o comentário, para o despertar de emoções fruto das percepções individuais, que conversam com as leituras do mundo.
Isso não significa atuar sem sistematização ou compromisso, mas conduzir para a construção de conhecimento a partir da diversão, satisfação e prazer.
Nesse percurso de construção do conhecimento é que se coloca a avaliação da compreensão leitora, auxiliando na tomada de decisões para a solução de problemas que promoverão a ampliação e a efetivação da aprendizagem.
A avaliação no eixo da produção textual deve ser pensada ressaltando o aprender a escrever relacionado ao uso de estratégias de produção de texto, adequando à situação sociocomunicativa.
Assim, a avaliação precisa ir ao encontro do objetivo principal : contribuir para a formação de um aluno autônomo produtor de textos, logo, numa perspectiva formativa e processual, auxiliando o aluno a amadurecer frente às dificuldades encontradas, de forma reflexiva e a partir da interlocução com situações reais de produção.
Nesse sentido, aparece como um recurso mediador, que irá oportunizar a adequação às condições de produção, contribuindo para o crescimento do produtor, a partir do atendimento às exigências da ordem do discurso, bem como da textualidade.
Isso se concretiza através da parceria professor/aluno na realização de um trabalho que envolve planejamento, escrita, revisão e reescrita, não, necessariamente, em uma ordem linear, mas como ações que acontecem de maneira concomitante : durante todo o processo de escrita há replanejamento, revisão e reescrita.
Por que avaliar leitura? > A avaliação como parte integrante do trabalho docente, deve, sim, incluir a \textbf{verificação} da capacidade de leitura do aluno.
Precisamos saber se o aluno compreende o que lê, porque isso é relevante para a vida em nossa sociedade letrada e porque, como professores, temos a responsabilidade de promover o desenvolvimento da competência leitora dos nossos alunos.
As representações do mundo manifestam -se em textos concretizados em diferentes gêneros, então, compreender textos é compreender o mundo, embora essa não seja a única maneira de fazê -lo ( BESERRA, 2007 p. 49 ). > A avaliação de redações só exercerá uma função formativa se, efetivamente, contribuir para que os alunos construam, consolidem e ampliem sua capacidade como escritores letrados, autônomos, críticos e historicamente situados ( MARCUSHI, 2007, p. 73 ). > A análise linguística não pode ter um fim em si mesma, tampouco pode ser avaliada isoladamente, pois corre -se o risco de fragmentar, de forma equivocada, o complexo fenômeno da linguagem ( MENDONÇA, 2007, p. 108 ).
A avaliação no eixo da oralidade também deve centrar -se na perspectiva formativa, reiterando que os gêneros orais da esfera escolar ( seminários, relatos de experiência, discussões em grupos, por exemplo ) e da esfera da vida pública ( testemunho, entrevista, negociação, debate, por exemplo ) devem estar presentes na sala de aula, porque conhecê -los e entender seu funcionamento é importante, já que eles estão muito relacionados ao exercício da cidadania.
Portanto, deve ultrapassar a reflexão sobre as variações dialetais e de registro ; deve prestigiar a análise de aspectos, tais como situações comunicativas e estratégias interacionais peculiares a cada gênero ; ou seja, é preciso atentar para o reconhecimento dos aspectos : extralinguísticos ( que remetem à relação entre os interlocutores ) ; paralinguísticos ( que remetem à fala propriamente dita, ao tom e ritmo de voz ), cinésicos ( relativos à expressão corporal ) ; e linguísticos ( que relacionam a construção do texto em si, como o uso de marcadores conversacionais, correções, hesitações, repetições, por exemplo ).
Na prática, isso quer dizer que o aluno deve ser avaliado quanto à percepção do uso de estratégias, pelo falante, em situações diversas de comunicação, para atender a necessidades também diversas.
A avaliação no eixo da análise linguística/semiótica deve ser associada à leitura e à produção textual, já que ela é parte integrante da construção dos sentidos dos textos.
Assim sendo, é importante que o professor oferte aos alunos situações que exijam a reflexão sobre a linguagem em seus próprios textos e em textos de outros.
Tal intervenção contribui para que os educandos amadureçam, melhor compreendam e produzam textos em gêneros diversos, fazendo escolhas de maneira consciente para atender aos objetivos do autor, revelando as nuances que perpassam as entrelinhas do texto, formando um leitor mais autônomo, no ambiente escolar e fora dele.
Na pretensão de avaliar a partir da compreensão do processo de aprendizagem, é imprescindível que o professor tenha em mente o seu papel como um mediador sempre disponível para a adaptação de suas práticas, em prol da ampliação da competência leitora e produtora de textos, visando ao desenvolvimento efetivo do aluno.
No contexto das etapas da alfabetização, a avaliação é concebida como uma forma de compreender os conhecimentos prévios das crianças para, a partir de então, planejar ações educativas que permitam diálogos construtivos entre saberes existentes e novos saberes, dessa forma, garantindo seus direitos de aprendizagem, sua intervenção na comunidade a qual pertencem, seu crescimento individual e coletivo.
Assim sendo, reitera -se que a avaliação é processual, acompanhada da observação das aprendizagens das crianças, compreendendo e respeitando os tempos de cada uma e a relevância de sua evolução dentro desses tempos, elaborando estratégias que as provoquem no processo de aprendizagem.
Por isso, não se avaliam os resultados ao final de cada etapa, mas todo processo de construção de conhecimentos vivenciados ao longo do percurso.
Essa compreensão de avaliação convida -nos a observar, refletir, planejar, intervir e, toda vez que se fizer necessário, ressignificar o processo.
Para todos esses momentos, é preciso espírito aberto e olhar sensível às diversas realidades, às particularidades e especificidades das crianças, que não chegam vazias de conhecimento e de sentido, mas trazem seus saberes, sua cultura, suas raízes, sua percepção de mundo, sua maneira de se relacionar com o outro e consigo mesmas, abrangendo aí toda a comunidade escolar.
Sendo assim, incluem -se na avaliação as ações docentes, a escola e o sistema de ensino.
Compreendido dessa forma, o processo avaliativo deve identificar, no final da etapa de alfabetização, se a criança faz uso autônomo da leitura/escrita, por exemplo, sem esquecer que a apropriação da escrita e a ampliação da oralidade devem favorecer a participação efetiva, em contextos particulares e públicos, nos mais distintos contextos da vida humana, articulando saberes e valores de sua comunidade.
Assim, a avaliação precisa ultrapassar as questões cognitivas e se abrir para a observação dos demais elementos inerentes ao processo de aquisição da escrita alfabética.
Por fim, a construção de um currículo engajado com a inclusão, como processos democráticos de avaliação, precisa romper com um formato avaliativo que prioriza resultados em detrimento dos processos formativos, estigmatizando os sujeitos desse processo.
Para muito além de tal perspectiva, a avaliação deve ser uma forma de ajudar, verdadeiramente, a incluir as crianças, as práticas docentes, a comunidade escolar como um todo, de forma ativa, reflexiva, participativa, efetiva, no processo educativo.

% matched lemmas: gradual, mensurar, verificação
\end{textsample}
